\begin{table}[htbp]
\centering
\small
\caption{Problemas B e C (openai/gpt-4o-mini): transferência zero-shot contra cada métrica estimada no A.}
\label{tab:T2_BC_gpt4mini}
\begin{tabular}{llllrrr}
\toprule
Prob. & Métrica de origem (A) & Algoritmo & Consistência (\%) & $\kappa$ & F1-Macro (\%) & Acur.\ LLM vs real \\
\midrule
B & zero-shot & nnls & 71,0 (68,0–76,0) & 0,358 (0,274–0,466) & 71,0 (68,0–76,0) & 57,7 (57,0–59,0) \\
B & zero-shot & perceptron & 70,7 (69,0–74,0) & 0,348 (0,288–0,421) & 70,7 (69,0–74,0) & 57,7 (57,0–59,0) \\
B & 4-shot & nnls & 71,4 (69,0–82,0) & 0,367 (0,288–0,599) & 71,4 (69,0–82,0) & 57,7 (57,0–59,0) \\
B & 4-shot & perceptron & 70,7 (67,0–79,0) & 0,354 (0,260–0,537) & 70,7 (67,0–79,0) & 57,7 (57,0–59,0) \\
B & 10-shot & nnls & 65,7 (48,0–83,0) & 0,259 (-0,070–0,589) & 65,7 (48,0–83,0) & 57,7 (57,0–59,0) \\
B & 10-shot & perceptron & 73,2 (67,0–83,0) & 0,407 (0,260–0,589) & 73,2 (67,0–83,0) & 57,7 (57,0–59,0) \\
\textbf{B} & \textbf{20-shot} & \textbf{nnls} & \textbf{64,8 (48,0–77,0)} & \textbf{0,250 (-0,070–0,483)} & \textbf{64,8 (48,0–77,0)} & \textbf{57,7 (57,0–59,0)} \\
\textbf{B} & \textbf{20-shot} & \textbf{perceptron} & \textbf{74,9 (67,0–83,0)} & \textbf{0,440 (0,260–0,589)} & \textbf{74,9 (67,0–83,0)} & \textbf{57,7 (57,0–59,0)} \\
B & 40-shot & nnls & 72,2 (53,0–85,0) & 0,395 (0,039–0,670) & 72,2 (53,0–85,0) & 57,7 (57,0–59,0) \\
B & 40-shot & perceptron & 77,1 (69,0–83,0) & 0,487 (0,334–0,591) & 77,1 (69,0–83,0) & 57,7 (57,0–59,0) \\
B & W ajustado ao rótulo real (teto) & nnls & 73,3 (69,0–77,0) & 0,409 (0,334–0,483) & 73,3 (69,0–77,0) & 57,7 (57,0–59,0) \\
B & W ajustado ao rótulo real (teto) & perceptron & 75,7 (69,0–83,0) & 0,457 (0,334–0,589) & 75,7 (69,0–83,0) & 57,7 (57,0–59,0) \\
\midrule
C & zero-shot & nnls & 88,3 (86,0–91,0) & 0,754 (0,699–0,812) & 88,3 (86,0–91,0) & 17,3 (16,0–19,0) \\
C & zero-shot & perceptron & 88,0 (85,0–90,0) & 0,745 (0,681–0,789) & 88,0 (85,0–90,0) & 17,3 (16,0–19,0) \\
C & 4-shot & nnls & 87,8 (84,0–90,0) & 0,743 (0,667–0,790) & 87,8 (84,0–90,0) & 17,3 (16,0–19,0) \\
C & 4-shot & perceptron & 88,7 (85,0–92,0) & 0,760 (0,681–0,830) & 88,7 (85,0–92,0) & 17,3 (16,0–19,0) \\
C & 10-shot & nnls & 84,8 (76,0–90,0) & 0,684 (0,501–0,790) & 84,8 (76,0–90,0) & 17,3 (16,0–19,0) \\
C & 10-shot & perceptron & 87,9 (85,0–92,0) & 0,745 (0,681–0,830) & 87,9 (85,0–92,0) & 17,3 (16,0–19,0) \\
\textbf{C} & \textbf{20-shot} & \textbf{nnls} & \textbf{85,4 (76,0–91,0)} & \textbf{0,698 (0,501–0,810)} & \textbf{85,4 (76,0–91,0)} & \textbf{17,3 (16,0–19,0)} \\
\textbf{C} & \textbf{20-shot} & \textbf{perceptron} & \textbf{87,6 (86,0–91,0)} & \textbf{0,739 (0,699–0,810)} & \textbf{87,6 (86,0–91,0)} & \textbf{17,3 (16,0–19,0)} \\
C & 40-shot & nnls & 86,3 (81,0–91,0) & 0,717 (0,613–0,810) & 86,3 (81,0–91,0) & 17,3 (16,0–19,0) \\
C & 40-shot & perceptron & 87,3 (85,0–92,0) & 0,735 (0,681–0,830) & 87,3 (85,0–92,0) & 17,3 (16,0–19,0) \\
C & W ajustado ao rótulo real (teto) & nnls & 87,0 & 0,732 (0,728–0,735) & 87,0 & 17,3 (16,0–19,0) \\
C & W ajustado ao rótulo real (teto) & perceptron & 87,0 (85,0–90,0) & 0,728 (0,681–0,790) & 87,0 (85,0–90,0) & 17,3 (16,0–19,0) \\
\bottomrule
\multicolumn{7}{p{0.95\linewidth}}{\footnotesize Rótulos ZERO-SHOT do LLM em B e C (100 pontos novos cada); só $\hat{W}$ é transferido do A (R$^2$), os centróides são recalculados com os rótulos do LLM no próprio problema (consistência). $\kappa$ = Cohen's $\kappa$ (faixas de Landis e Koch, 1977). F1-Macro entre a métrica e o LLM. Acur.\ LLM vs real = acurácia (vs rótulo real) do LLM, contexto. Em zero-shot (e com nomes de classe neutros) a correspondência entre nome e classe é arbitrária: a acurácia (vs rótulo real) do LLM pode sair abaixo de 50\%; o valor invariante à permutação é max($a$, 1 $-$ $a$). Em negrito: as linhas (Perceptron e NNLS) do regime/k eleito pelo runner (coluna melhor\_metrica\_a: maior fidelidade média no A, Perceptron, R$^2$, zero ou few-shot), excluído o teto.} \\
\end{tabular}
\end{table}
